\subsection{Spatial Occupancy Mapping \& Distance Fields}
Layered 2D costmaps \cite{lu2014layered, macenski2020marathon2} maintain geometric obstacles by ray-casting sensor ranges into grid cells and propagating configuration-space inflation. 3D volumetric frameworks, such as OctoMap \cite{hornung2013octomap}, employ hierarchical octrees with discrete probabilistic voxel occupancy updates. Distance-field frameworks such as Voxblox \cite{oleynikova2017voxblox} incrementally construct continuous Euclidean Signed Distance Fields (ESDF) for smooth trajectory optimization. These mapping representations are fundamentally designed to track geometric occupancy and distance fields from sensor rays. Because their core primitive is geometric spatial state rather than execution history, they do not inherently index episodic navigation action aborts or high-level failure metadata.

\subsection{Execution Monitoring \& Plan Diagnosis}
Deliberative robotics frameworks \cite{fikes1971strips, simmons1998xavier, ingrand2017deliberation} formalize plan execution monitoring and diagnostic repair. Modern navigation stacks, including Nav2 \cite{macenski2020marathon2}, execute local recovery behaviors (such as in-place rotations, costmap clearing, or backup maneuvers) during active trajectory tracking. In standard Nav2 configurations, navigation actions are dispatched independently, and recovery routines operate within the scope of an individual goal execution \cite{macenski2020marathon2}. Extending navigation to retain information across sequential goal dispatches can be addressed either through persistent global occupancy layers or via dedicated episodic failure stores at the decision layer.

\subsection{LLM Self-Reflection \& Agent Memory}
Recent embodied AI systems investigate large language models for planning and error correction. Reflexion \cite{shinn2023reflexion} prompts LLMs with scalar task feedback to generate natural language self-reflections across trials. Inner Monologue \cite{huang2023inner} integrates success detectors and scene descriptions into closed-loop prompt chains. REFLECT \cite{liu2023reflect} uses hierarchical multi-modal experience summarization with LLMs to generate structured failure explanations and plan corrections. While these frameworks demonstrate prompt-level flexibility over semantic descriptions, mobile robot navigation also requires grounding execution outcomes in metric sensor coordinates and spatial costmaps. FailMem addresses this complementary domain by investigating event-driven failure recording, perception-verified rule invalidation, and offline replay auditing directly within the ROS~2 Nav2 action stack.

\subsection{Experience-Based Planning \& No-Good Learning}
In motion planning, the Lightning framework \cite{berenson2012lightning} and Experience Graphs (E-Graphs) \cite{phillips2012egraphs} reuse libraries of prior successful paths to accelerate search. In discrete optimization, conflict-driven clause learning (e.g., GRASP \cite{marquessilva1999grasp}) records ``no-goods'' from search dead ends to prune infeasible branches. FailMem adapts the principle of negative constraint recording to mobile robot navigation with explicit sensor-verified invalidation.

Table~\ref{tab:classification} summarizes memory representations and state update/invalidation concepts across navigation and reasoning paradigms. Importantly, state update, ray clearing, graph invalidation, and symbolic rule expiration represent conceptually distinct mechanisms across paradigms: geometric mappers update cell occupancy via ray casting, path libraries perform collision checks during repair, LLM agents reset context across trials, while FailMem introduces perception-verified rule invalidation for recorded failure events.

\begin{table*}[t]
\centering
\small
\caption{Taxonomy of Memory Representations and State Update/Invalidation Concepts in Robot Navigation and Embodied Reasoning.}
\label{tab:classification}
\begin{tabular*}{\textwidth}{@{\extracolsep{\fill}}llll@{}}
\toprule
Paradigm & Representation & Grounding Level & State Update / Invalidation Concept \\
\midrule
\textbf{Costmap2D / Nav2} \cite{lu2014layered} & 2D Occupancy Grid & Geometric cells & Ray-clearing through free space \\
\textbf{OctoMap} \cite{hornung2013octomap} & 3D Voxel Octree & Discrete 3D voxels & Ray-casting probabilistic update \\
\textbf{Voxblox} \cite{oleynikova2017voxblox} & Continuous ESDF / TSDF & 3D Distance Field & Dynamic voxel weight integration \\
\textbf{Experience Graphs} \cite{phillips2012egraphs} & Path Waypoints / Graph & Discrete Path & Edge collision check during repair \\
\textbf{Reflexion} \cite{shinn2023reflexion} & Natural Language Text & Semantic tokens & Trial reset / verbal memory update \\
\textbf{REFLECT} \cite{liu2023reflect} & Natural Language Text & Multi-modal logs & Task execution cycle reset \\
\textbf{FailMem (Implemented)} & Event Tuple & Raw sensor bundle & Sensor-verified \texttt{FREE} classification \\
\bottomrule
\end{tabular*}
\end{table*}

